\subsection*{Merging paraphrased critiques}

The panel's value depends on recognising when two lenses raise the same problem in different words, both to avoid duplicate required fixes and to escalate issues on which lenses agree. Within the 162 panel instances, the primary judge's labels defined 1,342 cross-lens pairs of critiques identifying the same planted flaw and 3,214 pairs identifying different flaws. With word-level Jaccard similarity at a threshold set a priori without calibration ($\tau=0.5$), merging found 1.2\% of the duplicate pairs (pairwise F1 0.02; Fig~\ref{fig3}B): paraphrased critiques share too few exact words to reach that threshold. TF-IDF similarity at its a priori threshold ($\tau=0.3$) found 62\% (F1 0.76), because down-weighting common words lets paraphrases that share a few informative terms pass. The failure lies mostly in the threshold rather than in lexical similarity as such. With thresholds selected by leave-one-scenario-out cross-validation (0.075--0.1 in every fold; 0.075 in 11 of 12 folds for each method), lexical Jaccard reached a held-out F1 of 0.80 and TF-IDF cosine similarity 0.81, and LLM adjudication, which needs no threshold, reached 0.82 (Fig~\ref{fig3}A). On these flaw-matched pairs all three methods almost never merged critiques of different planted flaws (precision 0.997--1.000) and differed mainly in recall (0.67--0.70). This precision covers only the 29\% of panel critiques on flawed items that matched a planted flaw; it does not show that the remaining 71\%, unplanted concerns that carry no pair labels, were kept apart. The methods did differ in how much they merged overall: on flawed items the lexical methods formed 1,248 (Jaccard) and 1,259 (TF-IDF) multi-lens groups, against 858 for LLM adjudication, which left more critiques unmerged. Recall cannot exceed 0.74 in this design, because critiques from the same reviewer call are not merged under the platform default: when one lens raises a flaw twice and another lens raises it once, the single critique can join only one of the two, so one of the two cross-lens duplicate pairs cannot be recovered. In total, 349 of the 1,342 duplicate pairs are unreachable by construction; relative to that bound the methods recovered 91--95\% of the mergeable pairs. LLM adjudication returned a valid partition for all 162 panels at a cost of about 15,300 input and 1,300 output tokens per consensus step. The differences in F1 (0.01--0.02) are small, so we regard the three methods as comparable in duplicate detection; under the pre-specified selection rule (best held-out F1), LLM adjudication, the most conservative of the three, is the platform's default merge, with TF-IDF at $\tau=0.1$ as the zero-cost alternative. The optional sentence-embedding method, which requires a locally installed embedding model, was not included in this evaluation, and its default threshold is uncalibrated.

\textbf{Agreement is informative.} Merged groups raised by two or more lenses contained a planted flaw far more often than single-lens groups (49\% vs 5\% with LLM adjudication; 40\% vs 3\% with calibrated Jaccard), and 98\% of the planted flaws raised by two or more lenses were escalated. Escalation on agreement therefore concentrates severity on real problems (the escalation threshold itself, two distinct lenses, was not varied). It also has a side effect that the benchmark makes visible: when the benchmark grades were recomputed with TF-IDF merging at $\tau=0.075$ and escalation on agreement, every multi-call condition assigned grade F to every item, clean ones included, so the grade did not separate flawed from clean items (AUROC 0.50, against 0.66--0.76 for the underlying severity score; S2 Table). The grade discrimination reported above (AUROC 0.74--0.89 for the three-call conditions) therefore holds only when duplicates are largely left unmerged. TF-IDF stands in for LLM adjudication in this analysis because the adjudicator was run only on panel instances; the two methods had similar merge accuracy (F1 0.81 vs 0.82) and both escalated 98\% of the flaws raised by two or more lenses. Over-merging may contribute to the saturation: on clean items, TF-IDF merging formed more multi-lens groups than LLM adjudication (524 vs 308 over the 54 clean panels), and the ground truth cannot show whether such groups join one concern or several. We return to the consequences for grading below.
